\exercisesection{\S~3}

\begin{enumerate}
\def\labelenumi{\arabic{enumi})}
\item
  设 \(\Gamma\) 是 \(\mathrm{Q}(\mathrm{R})\) 的有限指数子群，且 \(\mathrm{P} = \Gamma \cap \mathrm{R}\)。代数 \(\mathfrak{h} + \mathfrak{g}^{\mathrm{P}}\) 是约化的，并且 \(\mathfrak{g}\) 中任何包含 \(\mathfrak{h}\) 的约化子代数都可以用这种方式得到（利用第 VI 章，§1，习题 6 b)）。
\item
  令 X 为 g 的约化子代数集合中与 g 不同且包含 h 的那些。利用第 VI 章，§4，习题 4 确定 X 的极小元素。证明这种极大子代数的中心维数为 0 或 1，分别对应于所讨论习题的 a) 或 b) 情形。
\item
  设 \(\mathfrak{b}\) 是 \(\mathfrak{g}\) 的一个 Borel 子代数，且 \(l = \mathrm{rk}(\mathfrak{g})\)。证明代数 \(\mathfrak{b}\) 的最少生成元数是 \(l\)，当 \(l \neq 1\) 时，且在 \(l = 1\) 时为 2。
\item
  假设 \(k = \mathbf{R}\) 或 \(\mathbf{C}\)。置 \(G = \operatorname{Int}(\mathfrak{g})\)，并将 \(G\) 的李代数与 \(\mathfrak{g}\) 认同。设 \(\mathfrak{b}\) 是 \((\mathfrak{g},\mathfrak{h})\) 的一个 Borel 子代数，且 \(\mathfrak{n}\) 是其幂零元素的集合。记 \(H,B,N\) 为 \(G\) 中李代数分别为 \(\mathfrak{h},\mathfrak{b},\mathfrak{n}\) 的积分子群。证明 \(H,B,N\) 是 \(G\) 的李子群，\(N\) 是单连通的，且 \(B\) 是 \(H\) 与 \(N\) 的半直积。
\item
  设 \(\mathfrak{m}\) 是半单李代数 \(\mathfrak{a}\) 的一个抛物子代数。
\end{enumerate}

\begin{enumerate}
\def\labelenumi{\alph{enumi})}
\item
  设 p 是 m 的子代数。证明 p 是 a 的抛物子代数当且仅当 p 包含 m 的根基 r，且 p/r 是半单代数 m/r 的抛物子代数。
\item
  若 \(m'\) 是 a 的一个抛物子代数，则 \(m \cap m'\) 的每个嘉当子代数都是 a 的嘉当子代数。（归约到分裂情形，并将命题 10 应用于包含在 m 和 \(m'\) 中的 Borel 子代数。）
\end{enumerate}

\begin{enumerate}
\def\labelenumi{\arabic{enumi})}
\setcounter{enumi}{5}
\item
  半单李代数 \(\mathfrak{a}\) 的两个 Borel 子代数若其交是嘉当子代数，则称它们相反。证明，如果 \(\mathfrak{b}\) 是 \(\mathfrak{a}\) 的一个 Borel 子代数，且 \(\mathfrak{h}\) 是 \(\mathfrak{b}\) 的一个嘉当子代数，则存在唯一的 \(\mathfrak{a}\) 的 Borel 子代数，它与 \(\mathfrak{b}\) 相反且包含 \(\mathfrak{h}\)。（归约到分裂情形。）
\item
  设 \(\mathfrak{a}\) 是半单李代数，\(\mathfrak{h}\) 是 \(\mathfrak{a}\) 的嘉当子代数，且 \(\mathfrak{s}\) 是 \(\mathfrak{a}\) 的一个包含 \(\mathfrak{h}\) 的半单子代数。证明：
\end{enumerate}

\[
(\mathfrak {a}, \mathfrak {h}) \text {is split} \Longleftrightarrow (\mathfrak {s}, \mathfrak {h}) \text {is split.}
\]

构造一个例子，其中 \(\mathfrak{a}\) 可分裂但 \(\mathfrak{s}\) 不可分裂。（取 \(\mathfrak{a} = \mathfrak{sp}(4, k)\) 和 \(\mathfrak{s} = \mathfrak{sl}(2, k')\)，其中 \(k'\) 是 \(k\) 的二次扩域。）

\begin{enumerate}
\def\labelenumi{\arabic{enumi})}
\setcounter{enumi}{7}
\tightlist
\item
  选择 R 的一个基，因而得到一组正根 \(R_{+}\)。假设 R 不可约。令 \(\tilde{\alpha}\) 为最高根，S 为与 \(\tilde{\alpha}\) 正交的根的集合，且 \(S_{+} = S \cap R_{+}\)。
\end{enumerate}

\begin{enumerate}
\def\labelenumi{\alph{enumi})}
\tightlist
\item
  令
\end{enumerate}

\[
\mathfrak {g} ^ {\prime} = \mathfrak {g} ^ {\mathrm{S}} \oplus \mathfrak {h} _ {\mathrm{S}}, \quad \mathfrak {m} _ {+} = \sum_ {\alpha \in \mathrm{S} _ {+}} \mathfrak {g} ^ {\alpha}, \quad \mathfrak {m} _ {-} = \sum_ {\alpha \in \mathrm{S} _ {+}} \mathfrak {g} ^ {- \alpha}.
\]

则 \(\mathfrak{g}'\) 是半单的且 \(\mathfrak{g}' = \mathfrak{h}_{\mathrm{S}} \oplus \mathfrak{m}_{+} \oplus \mathfrak{m}_{-}\)。

\begin{enumerate}
\def\labelenumi{\alph{enumi})}
\setcounter{enumi}{1}
\tightlist
\item
  置：
\end{enumerate}

\[
\mathfrak {n} _ {+} = \sum_ {\alpha \in \mathrm{R} _ {+}} \mathfrak {g} ^ {\alpha}, \quad \mathfrak {p} = \sum_ {\alpha \in \mathrm{R} _ {+} - \mathrm{S} _ {+}} \mathfrak {g} ^ {\alpha}, \quad \mathfrak {p} _ {0} = \sum_ {\alpha \in \mathrm{R} _ {+} - \mathrm{S} _ {+} - \{\tilde {\alpha} \}} \mathfrak {g} ^ {\alpha}.
\]

则

\[
\mathfrak {n} _ {+} = \mathfrak {m} _ {+} \oplus \mathfrak {p}, \quad [ \mathfrak {m} _ {+}, \mathfrak {p} _ {0} ] \subset \mathfrak {p} _ {0}, \quad [ \mathfrak {n} _ {+}, \mathfrak {g} ^ {\tilde {\alpha}} ] = 0.
\]

特别地，\(\mathfrak{p}\) 是 \(\mathfrak{n}_{+}\) 的理想。

\begin{enumerate}
\def\labelenumi{\alph{enumi})}
\setcounter{enumi}{2}
\tightlist
\item
  对所有 \(\alpha \in \mathrm{R}_{+} - \mathrm{S}_{+} - \{\tilde{\alpha}\}\)，存在唯一的 \(\alpha' \in \mathrm{R}_{+} - \mathrm{S}_{+} - \{\tilde{\alpha}\}\)，使得 \(\alpha + \alpha' = \tilde{\alpha}\)。对所有 \(\alpha \in \mathrm{R}_{+} - \mathrm{S}_{+} - \{\tilde{\alpha}\}\)，可以选择非零元素 \(X_{\alpha}\)，它属于 \(\mathfrak{g}^{\alpha}\)，使得
\end{enumerate}

\[
\begin{array}{r l} {[ X _ {\alpha}, X _ {\alpha^ {\prime}} ] = \pm X _ {\tilde {\alpha}}} & {\mathrm{if} \alpha \in \mathrm{R} _ {+} - \mathrm{S} _ {+} - \{\tilde {\alpha} \}} \\ {[ X _ {\alpha}, X _ {\beta} ] = 0} & {\mathrm{if} \alpha , \beta \in \mathrm{R} _ {+} - \mathrm{S} _ {+}, \beta \neq \alpha^ {\prime}. ^ {1 1}} \end{array}
\]

\footnotetext[11]{本习题由 A. JOSEPH 提供给我们。}

\begin{enumerate}
\def\labelenumi{\arabic{enumi})}
\setcounter{enumi}{8}
\tightlist
\item
  构造半单李代数 \(\mathfrak{a}\) 的例子，使得：
\end{enumerate}

\begin{enumerate}
\def\labelenumi{\roman{enumi})}
\item
  a 没有 Borel 子代数；
\item
  a 有一个 Borel 子代数，但不是可分裂的。
\end{enumerate}

¶10) 设 a 是半单李代数，且 x 是 a 的一个元素。如果 ad x 可对角化（《代数》第 VII 章），换言之，如果存在一个由 ad x 的特征向量组成的 a 的基，则称 x 是可对角化的。

\begin{enumerate}
\def\labelenumi{\alph{enumi})}
\item
  设 \(\mathfrak{c}\) 是 \(\mathfrak{a}\) 的一个交换子代数，由可对角化元素组成，且 L 是 \(\mathfrak{c}\) 在表示 \(\mathrm{ad}_{\mathfrak{a}}\) 中的权集。集合 L 是 \(\mathfrak{c}^*\) 的一个有限子集，包含 0（除非 \(\mathfrak{a} = 0\)），且 \(\mathfrak{a} = \bigoplus_{\lambda \in L} \mathfrak{a}^{\lambda}(\mathfrak{c})\)。证明存在 L 的子集 M，使得 \(L - \{0\}\) 是 M 与 \(-M\) 的不交并，且 \((M + M) \cap L \subset M\)。如果 M 具有这些性质，置 \(\mathfrak{a}^{M} = \bigoplus_{\lambda \in M} \mathfrak{a}^{\lambda}(\mathfrak{c})\)，并置 \(\mathfrak{p}^{M} = \mathfrak{a}^{0}(\mathfrak{c}) \oplus \mathfrak{a}^{M}\)。代数 \(\mathfrak{a}^{0}(\mathfrak{c})\) 是 \(\mathfrak{c}\) 在 \(\mathfrak{a}\) 中的交换子；它在 \(\mathfrak{a}\) 中是约化的（第 VII 章，§1，第 5 节），且其嘉当子代数就是 \(\mathfrak{a}\) 的嘉当子代数（第 VII 章，§2，第 3 节）。证明 \(\mathfrak{p}^{M}\) 是 a 的抛物子代数，且 \(a^{M}\) 是 \(p^{M}\) 的根基的幂零元素集合。（利用 c 包含于 a 的一个嘉当子代数这一事实，并通过扩张标量归约到这个子代数是分裂的情形。）
\item
  保留 a) 的假设和记号，并进一步假设 a 可分裂。证明 c 包含于 a 的一个分裂嘉当子代数中。（取 a 的一个包含于 \(p^{M}\) 中的分裂嘉当子代数 h；则存在唯一的 \(h'\)，它是 \(\mathfrak{a}^{0}(\mathfrak{c})\) 的嘉当子代数，使 h 包含于 \(h' \oplus a^{M}\) 中；证明 \(h'\) 是 a 的分裂嘉当子代数，且 \(h'\) 包含 c。）
\end{enumerate}

\begin{enumerate}
\def\labelenumi{\arabic{enumi})}
\setcounter{enumi}{10}
\tightlist
\item
  设 \(\mathfrak{a} = \prod \mathfrak{a}_i\) 是半单李代数的有限积。\(\mathfrak{q}\) 是 \(\mathfrak{a}\) 的子代数，且它是抛物（分别为 Borel）子代数，当且仅当它形如 \(\mathfrak{q} = \prod \mathfrak{q}_i\)，其中对于所有 \(i\)，\(\mathfrak{q}_i\) 是 \(\mathfrak{a}_i\) 的抛物（分别为 Borel）子代数。
\end{enumerate}

¶ 12) 设 \(k'\) 是 \(k\) 的有限扩域，\(\mathfrak{a}'\) 是半单 \(k'\)-李代数，且 \(\mathfrak{a}\) 是其底层 \(k\)-李代数。证明 \(\mathfrak{a}\) 的抛物（分别为 Borel）子代数与 \(\mathfrak{a}'\) 的相同。（将标量扩张到 \(k\) 的代数闭包，并利用习题 11。）

\begin{enumerate}
\def\labelenumi{\arabic{enumi})}
\setcounter{enumi}{12}
\item
  设 p 和 q 是半单李代数 a 的两个抛物子代数，且 n 是 p 的幂零根基。证明 \(\mathfrak{m} = (\mathfrak{p} \cap \mathfrak{q}) + \mathfrak{n}\) 是 a 的抛物子代数。（归约到 a 分裂且 q 为 Borel 子代数的情形；选择包含于 \(p \cap q\) 中的嘉当子代数 h，并确定相应根系的子集 P，使 \(m = h + g^{P}\)。）
\item
保留命题 9 的记号。
\end{enumerate}

\begin{enumerate}
\def\labelenumi{\alph{enumi})}
\item
  令 \(\alpha \in \mathrm{B}\)。证明 \(\mathfrak{n} \cap \operatorname{Kerad} X_{\alpha}\) 是各 \(\mathfrak{g}^{\beta}\) 的直和，其中 \(\beta\) 属于 \(\mathrm{R}_{+}\) 中满足 \(\alpha + \beta \notin \mathrm{R}_{+}\) 的元素集合。
\item
  推出，如果 g 是单的，则 n 的中心等于 \(g^{\tilde{\alpha}}\)，其中 \(\tilde{\alpha}\) 是 \(R_{+}\) 的最高根。在一般情形下，n 的中心维数等于 g 的单分量数。
\end{enumerate}
% semantic-fragments: legacy-input-seam
\relax{}
